\chapter{Concluzii}\label{ch:concluzii}
\pagestyle{fancy}

%==========================================================================
\section{Contribuții}
\label{sec:contributii}

Prezenta lucrare a realizat proiectarea și implementarea completă a unui
sistem de control FPGA pentru o mână robotică cu 8~servomotoare, integrând
filtrare Kalman pentru stabilizarea mișcării. Principalele contribuții sunt:

\begin{enumerate}
	\item \textbf{Control în timp real pe FPGA cu filtrare Kalman.}
	      S-a implementat un filtru Kalman cu 2~stări (poziție și viteză) în
	      aritmetică virgulă fixă Q16.16, multiplexat în timp pe 8~canale.
	      Filtrul estimează optimal traiectoria fiecărui servomotor pe baza
	      măsurătorilor zgomotoase, reducând fluctuațiile cu un factor de
	      4--6$\times$ comparativ cu controlul direct. Întregul ciclu de
	      filtrare pentru cele 8~canale se execută în aproximativ 1\,456
	      cicluri de ceas (14,6\,$\mu$s), reprezentând doar 0,07\% din
	      bugetul temporal disponibil de 20\,ms.

	\item \textbf{Generare PWM paralelă pe 8~canale.} Cele 8~semnale PWM
	      sunt generate simultan de instanțe hardware independente, fără
	      multiplexare în timp. Aceasta elimină jitter-ul de interleaving
	      specific implementărilor bazate pe microcontroler, rezultând o
	      precizie temporală de 10\,ns (un ciclu de ceas) pe fiecare canal.

	\item \textbf{Integrarea feedback-ului ADC.} Controlerul AD7124-8
	      citește poziția reală a celor 8~servomotoare prin potențiometrele
	      de feedback, furnizând măsurători pe 24~de biți la o rată de
	      50\,Hz. Aceste măsurători alimentează filtrul Kalman, permițând
	      estimarea nu doar a poziției, ci și a vitezei de deplasare.

	\item \textbf{Mod demonstrativ autonom.} ROM-ul cu 21~de cadre-cheie
	      permite funcționarea autonomă a sistemului, fără conexiune la un
	      calculator gazdă, facilitând demonstrarea și testarea rapidă.

	\item \textbf{Arhitectură modulară completă.} Cele 12~module VHDL sunt
	      proiectate cu interfețe clar definite și pot fi reutilizate
	      independent în alte proiecte. Gradul scăzut de utilizare a
	      resurselor FPGA (sub 15\%) lasă marjă substanțială pentru
	      extinderi ulterioare.
\end{enumerate}

%==========================================================================
\section{Analiză critică}
\label{sec:analiza-critica}

\subsection{Puncte forte}

\begin{itemize}
	\item \textbf{Latență redusă:} Latența de procesare a FPGA-ului
	      (de la recepția cadrului UART la actualizarea semnalelor PWM) este
	      sub 15\,$\mu$s, cu trei ordine de mărime mai mică decât ciclul
	      servo de 20\,ms. Aceasta garantează că procesarea nu introduce
	      întârzieri perceptibile.

	\item \textbf{Paralelism real:} Generarea simultană a tuturor canalelor
	      PWM, funcționarea concurentă a receptorului UART, a controlerului
	      SPI și a filtrului Kalman --- toate în același ciclu de ceas ---
	      demonstrează avantajul fundamental al FPGA față de un procesor
	      secvențial.

	\item \textbf{Utilizare eficientă a resurselor:} Sistemul complet
	      ocupă sub 15\% din resursele FPGA-ului Artix-7 XC7A35T, cel mai
	      mic circuit din familia Artix-7. Aceasta demonstrează că
	      implementarea este eficientă din punct de vedere al resurselor
	      hardware.

	\item \textbf{Robustețe la comunicare:} Protocolul de cadre cu
	      checksum XOR, timeout de 1\,ms și buffer dublu (mailbox) asigură
	      funcționarea stabilă chiar și în condiții de comunicare UART
	      imperfectă.
\end{itemize}

\subsection{Limitări}

\begin{itemize}
	\item \textbf{Filtru Kalman cu 2 stări:} Modelul curent nu include
	      accelerația ca variabilă de stare. Pentru mișcări rapide cu
	      accelerație variabilă, estimarea poate rămâne în urma mișcării
	      reale. Un filtru cu 3~stări (poziție, viteză, accelerație) ar
	      captura mai fidel dinamica mișcărilor brusteri, dar necesită
	      o matrice de covarianță $3\times3$ (6~elemente unice) în loc de
	      $2\times2$ (3~elemente).

	\item \textbf{Constante de ajustare fixe:} Constantele filtrului Kalman
	      ($Q_{00}$, $Q_{11}$, $R$) sunt hardcodate în VHDL și necesită
	      resinteză pentru modificare. Valorile actuale sunt calculate
	      teoretic și ar beneficia de calibrare pe hardware-ul real,
	      eventual cu un mecanism de ajustare în timp real prin UART.

	\item \textbf{Comunicare doar prin cablu:} Sistemul depinde de o
	      conexiune USB cablată pentru comunicarea UART cu calculatorul
	      gazdă. Aceasta limitează mobilitatea operatorului și introduce
	      o dependență fizică de proximitatea calculatorului.

	\item \textbf{Configurare ADC fără validare:} Secvența de inițializare
	      a AD7124-8 scrie registrele de configurare fără a citi înapoi
	      valorile pentru confirmare. O eroare de comunicare SPI în timpul
	      inițializării ar putea lăsa ADC-ul într-o stare incorectă fără
	      detecție.
\end{itemize}

%==========================================================================
\section{Dezvoltări ulterioare}
\label{sec:dezvoltari}

Arhitectura modulară și gradul scăzut de utilizare a resurselor FPGA
permit mai multe direcții de extindere:

\begin{enumerate}
	\item \textbf{Filtru Kalman cu 3 stări.} Adăugarea accelerației ca a
	      treia variabilă de stare ar îmbunătăți urmărirea mișcărilor
	      rapide. Vectorul de stare devine $[x_0, x_1, x_2]$ (poziție,
	      viteză, accelerație), iar matricea de covarianță crește de la
	      $2\times2$ la $3\times3$. Termenul $dt^2/2$ (0,0002 în Q16.16)
	      necesită o reprezentare cu mai mulți biți fracționari (Q4.28)
	      pentru a evita pierderea de precizie. Estimarea de resurse indică
	      45~multiplicări și 3~diviziuni per canal, încadrabile confortabil
	      în bugetul temporal existent.

	\item \textbf{Comunicare wireless.} Înlocuirea legăturii USB-UART cu
	      un modul Bluetooth (ex. HC-05) sau WiFi (ex. ESP32) ar permite
	      operarea la distanță. Modulele \texttt{uart\_rx} și
	      \texttt{uart\_tx} ar rămâne nemodificate, deoarece aceste module
	      wireless comunică prin UART standard. Singura modificare ar fi
	      la nivel de placă, conectând modulul wireless la pinii RsRx/RsTx
	      în locul convertorului USB-UART integrat.

	\item \textbf{Buclă de control PID.} Combinarea ieșirii filtrului
	      Kalman (estimarea poziției reale) cu unghiul țintă (din cadrul
	      UART) într-un controler PID ar permite compensarea activă a
	      erorilor de poziționare ale servomotoarelor. Eroarea
	      $e = \theta_{\text{tinta}} - \theta_{\text{Kalman}}$ ar alimenta
	      un controler PID a cărui ieșire ar ajusta comanda PWM, închizând
	      bucla de reglare.

	\item \textbf{Clasificare gesturi prin învățare automată.} Integrarea
	      unui clasificator simplu (rețea neuronală compactă sau arbore
	      de decizie) implementat în FPGA ar permite recunoașterea gesturilor
	      direct pe placă, fără a depinde de procesarea MediaPipe pe un
	      calculator gazdă. Gesturile recunoscute ar putea declanșa secvențe
	      predefinite, similar modului demo actual, dar bazate pe feedback-ul
	      ADC în timp real.

	\item \textbf{Extinderea numărului de grade de libertate.} Arhitectura
	      actuală suportă 8~canale, dar poate fi extinsă la 16 sau mai
	      multe prin adăugarea de instanțe PWM și extinderea filtrului
	      Kalman. Resursele FPGA disponibile (peste 85\% neutilizate)
	      permit dublarea numărului de canale fără dificultăți.
\end{enumerate}
